\chapter{Introduction}
\section{Giới thiệu đề tài và Đặt vấn đề}

\subsection{Lý do chọn đề tài (Motivation)}

Trong kỷ nguyên chuyển đổi số, nhu cầu về việc quản trị và khai thác tri thức đã trở nên cấp thiết hơn bao giờ hết. Các hệ quản trị cơ sở dữ liệu truyền thống (RDBMS) mặc dù rất mạnh mẽ trong việc lưu trữ và truy xuất dữ liệu có cấu trúc, nhưng vẫn còn nhiều hạn chế trong việc tự động suy diễn ra các tri thức mới từ tập dữ liệu hiện có. 

Việc tích hợp trí tuệ nhân tạo, cụ thể là các hệ thống dựa trên tri thức (Knowledge-based systems), vào tầng lưu trữ dữ liệu giúp rút ngắn khoảng cách giữa "Dữ liệu thô" và "Tri thức hữu ích" [1]. Đề tài "Xây dựng Hệ quản trị Cơ sở tri thức (KBMS) dựa trên mô hình Đối tượng tính toán" được lựa chọn nhằm giải quyết bài toán này.

\subsection{Mục tiêu nghiên cứu (Research Objectives)}

Mục tiêu chính của đề tài là xây dựng một hệ thống KBMS hoàn chỉnh, có khả năng:
1.  \textbf{Biểu diễn tri thức chuyên sâu}: Sử dụng mô hình COKB [1] để định nghĩa các thực thể có khả năng tính toán.
2.  \textbf{Suy diễn tự động hiệu năng cao}: Áp dụng thuật toán F-Closure để tìm tập đóng tri thức một cách tối ưu.
3.  \textbf{Lưu trữ bền vững và an toàn}: Phát triển công cụ lưu trữ nhị phân hỗ trợ chỉ mục B+ Tree và nhật ký WAL [5].
4.  \textbf{Giao diện phát triển trực quan}: Cung cấp Studio IDE giúp người dùng dễ dàng thiết kế và kiểm chứng tri thức.

\subsection{Đối tượng và Phạm vi nghiên cứu}

*   \textbf{Đối tượng}: Các mô hình biểu diễn tri thức, thuật toán suy diễn tiến (Forward Chaining) và các kỹ thuật quản lý CSDL.
*   \textbf{Phạm vi}: 
    *   \textbf{Quản lý tri thức}: Thực hiện các thao tác thêm, sửa, xóa và tìm kiếm tri thức chuyên sâu qua ngôn ngữ KBQL.
    *   \textbf{Suy diễn tri thức}: Xây dựng bộ máy suy duyễn dựa trên tập luật và sự thật hiện có để sinh ra tri thức mới.
    *   \textbf{Lưu trữ}: Nghiên cứu phương pháp lưu trữ tri thức bền vững dưới cấu trúc mô hình đối tượng tính toán (COKB).

\subsection{Ý nghĩa khoa học và thực tiễn}

Đề tài góp phần chuẩn hóa phương pháp xây dựng hệ thống tri thức có khả năng mở rộng (Scale-out) và tính linh hoạt cao trong việc thay đổi luật mà không cần can thiệp vào mã nguồn ứng dụng. Đây là nền tảng quan trọng cho việc phát triển các hệ chuyên gia và hệ trợ giúp quyết định thông minh.

